\documentclass[11pt,a4paper]{article}

%%% ── packages ────────────────────────────────────────────────────────
\usepackage[utf8]{inputenc}
\usepackage[T1]{fontenc}
\usepackage{amsmath,amssymb,amsthm}
\usepackage{mathtools}
\usepackage{booktabs}
\usepackage[dvipsnames]{xcolor}
\usepackage{hyperref}
\usepackage{tcolorbox}

%%% ── colour commands ─────────────────────────────────────────────────
\newcommand{\old}[1]{{\color{Red}#1}}          % flawed / old approach
\newcommand{\new}[1]{{\color{Blue}#1}}         % corrected approach

%%% ── coloured boxes for contrasting old vs new ───────────────────────
\newtcolorbox{oldbox}{colframe=red!70!black, colback=red!5, title={\textbf{Old}}}
\newtcolorbox{newboxelem}{colframe=blue!70!black, colback=blue!5, title={\textbf{New}}}

%%% ── theorem-like environments ───────────────────────────────────────
\newtheorem{definition}{Definition}
\newtheorem{problem}{Problem}
\newtheorem{lemma}{Lemma}
\newtheorem{theorem}{Theorem}
\newtheorem{proposition}{Proposition}
\newtheorem{corollary}{Corollary}

%%% ── macros ──────────────────────────────────────────────────────────
\newcommand{\A}{\mathcal{A}}
\newcommand{\loc}{\mathrm{loc}}
\newcommand{\act}{\mathrm{act}}
\newcommand{\Nodes}{N}
\newcommand{\Net}{\mathcal{G}}
\newcommand{\T}{\mathcal{T}}
\newcommand{\Tl}{\T^{\mathrm{loc}}}
\newcommand{\Tc}{\T^{\mathrm{c}}}
\newcommand{\Lg}{\mathcal{L}}
\newcommand{\cost}{\mathrm{cost}}
\newcommand{\alignment}{\gamma}
\newcommand{\SKIP}{\mathord{\gg}}
\newcommand{\trace}{\sigma}
\newcommand{\na}[2]{#1(#2)}
\newcommand{\depth}{\mathrm{depth}}
\newcommand{\ep}{\mathrm{entryPts}}
\newcommand{\ls}{\mathrm{localStarts}}
\newcommand{\GA}{\textsc{GetAlign}}
\newcommand{\cov}{\mathrm{covered}}
\newcommand{\visiting}{\mathrm{visiting}}

\title{Formal Problem Definition\\[4pt]
\large Distributed Online Conformance Checking via\\
Trie-Based Alignment Decomposition}
\author{}
\date{}

\begin{document}
\maketitle

%% =====================================================================
\section{Definitions}
\label{sec:definitions}
%% =====================================================================

We begin by establishing all definitions required for the formal problem
statement and the subsequent proof of soundness.

% ------------------------------------------------------------------
\subsection{Located Activities and Traces}
% ------------------------------------------------------------------

\begin{definition}[Located activity]\label{def:located-activity}
  Let $\A$ be a finite activity alphabet and
  $\Nodes = \{n_1, \dots, n_k\}$ a finite set of network nodes.
  A \emph{located activity} is a pair
  $\ell = (a, n) \in \A \times \Nodes$, written $\na{n}{a}$.
  We define $\act(\ell) = a$ and $\loc(\ell) = n$.
\end{definition}

\begin{oldbox}
\textbf{Flaw: Activity naming with occurrence counters.}\;
Currently, we construct activity labels as
$a' = \texttt{name\text{-}location\text{-}occurrence}$
(e.g.\ \texttt{CRP-B0}, \texttt{CRP-B1}), embedding the occurrence
index into the activity name.
Formally, this replaced $\A$ with a much larger alphabet
$\A' = \A \times \Nodes \times \mathbb{N}$, making the trie
hyper-specific to training data and destroying generalisation.
\end{oldbox}

\begin{newboxelem}
\textbf{Fix: Canonical activity names.}\;
Activities use the original name only: $a \in \A$.
The located activity $(a, n)$ separates identity from location.
\end{newboxelem}

\medskip

\begin{definition}[Located activity equality]\label{def:la-equality}
  Two located activities $\ell_1 = (a_1, n_1)$ and $\ell_2 = (a_2, n_2)$
  are \emph{equal}, written $\ell_1 = \ell_2$, if and only if
  $a_1 = a_2$ \textbf{and} $n_1 = n_2$.

  A weaker predicate $\mathrm{sameActivity}(\ell_1, \ell_2)$ holds iff
  $a_1 = a_2$ (ignoring location).
\end{definition}

\begin{oldbox}
\textbf{Flaw: Location-ignoring equality.}\;
The original equality checked only $a_1 = a_2$, ignoring the location
component.  This caused incorrect trie matches: an activity at node
$n_1$ would match the same activity name at node $n_2$, producing
wrong alignments.
\end{oldbox}

\begin{newboxelem}
\textbf{Fix.}\;
Equality requires both $\act(\ell_1) = \act(\ell_2)$ and
$\loc(\ell_1) = \loc(\ell_2)$.
The hash function is $h(\ell) = \mathrm{hash}(a, n)$.
\end{newboxelem}

\begin{definition}[Located trace]\label{def:located-trace}
  A \emph{located trace} is a finite, temporally ordered sequence
  \[
    \trace = \bigl\langle
      \ell_1^{t_1},\;
      \ell_2^{t_2},\;
      \ldots,\;
      \ell_m^{t_m}
    \bigr\rangle,
    \qquad
    \ell_i = (a_i, n_{j_i}) \in \A \times \Nodes,
    \quad
    t_1 < t_2 < \cdots < t_m.
  \]
  The \emph{local projection} $\trace|_{n}$ retains only events
  whose location is~$n$.
\end{definition}

\begin{definition}[Event log]\label{def:event-log}
  An \emph{event log} is a finite multiset
  $L = \{\trace_1, \trace_2, \ldots, \trace_p\}$
  of located traces.
\end{definition}

% ------------------------------------------------------------------
\subsection{Process Model}
% ------------------------------------------------------------------

\begin{definition}[Trie model]\label{def:trie}
  A \emph{trie model} is a rooted, labelled tree
  $\T = (Q, q_0, \delta, \lambda, F)$
  where $Q$ is a finite state set, $q_0 \in Q$ the root,
  $\delta \subseteq Q \times Q$ the parent--child relation,
  $\lambda\colon Q \setminus \{q_0\} \to \A \times \Nodes$ labels each
  non-root state with a located activity, and
  $F \subseteq Q$ is the set of \emph{accepting states}.

  Every root-to-accepting-state path encodes one admissible
  located trace.  We write $\Lg(\T)$ for the \emph{language} of~$\T$.
\end{definition}

\begin{definition}[Trie label reachability]\label{def:reachability}
  For a subtrie rooted at $q \in Q$, the \emph{reachable label set} is
  \[
    R(q) = \{ \lambda(q') \mid q' \in Q,\;
              q' \text{ is a descendant of } q \} \cup \{ \lambda(q) \}.
  \]
  We say a node $q$ \emph{contains} label~$\ell$ iff $\ell \in R(q)$.
\end{definition}

\begin{oldbox}
\textbf{Flaw: Recursive reachability check.}\;
Computed $R(q)$ via deep recursion:
each call to \texttt{contains($\ell$)} triggered $O(|Q|)$ recursive
descent, causing stack overflow on large tries.
\end{oldbox}

\begin{newboxelem}
\textbf{Fix: Iterative BFS with caching.}\;
$R(q)$ is built once via iterative breadth-first traversal and cached.
Subsequent containment queries are $O(1)$ set lookups.
Children are additionally indexed by label in a hash map for $O(1)$
child access.
\end{newboxelem}

\begin{definition}[Distributed trie model]\label{def:dist-trie}
  Given $\T$ over $\Nodes$, the \emph{distributed decomposition} assigns
  to each $n_i \in \Nodes$ a \emph{local trie} $\Tl_{n_i}$.
  Every maximal contiguous subsequence of activities at~$n_i$ in any
  root-to-accepting-state path of~$\T$ forms a branch in~$\Tl_{n_i}$.

  When a path transitions from $n_j$ to $n_i$
  ($\na{n_j}{a'} \to \na{n_i}{a}$, $n_j \neq n_i$),
  $\na{n_j}{a'}$ is inserted as an \emph{entry point} at the root
  of~$\Tl_{n_i}$; the subtrie below it contains the local continuation
  at~$n_i$.

  Formally,
  $\Tl_{n_i} = (Q_{n_i}, q_0^{n_i}, \delta_{n_i}, \lambda_{n_i}, F_{n_i})$
  with root children partitioned as:
  \begin{align*}
    \ls(n_i) &= \{q \mid \lambda_{n_i}(q)=\na{n_i}{a},\;
      \text{$a$ starts a trace at $n_i$}\},\\
    \ep(n_i) &= \{q \mid \lambda_{n_i}(q)=\na{n_j}{a'},\;
      n_j \neq n_i,\;
      \text{$\na{n_j}{a'}$ precedes a transition to $n_i$}\},\\
    \mathrm{children}(q_0^{n_i})
      &= \ls(n_i) \cup \ep(n_i).
  \end{align*}
\end{definition}

\begin{definition}[Network]\label{def:network}
  A \emph{network}
  $\Net = (\Nodes, \{\Tl_{n_i}\}_{n_i \in \Nodes})$
  pairs a node set with its local tries from
  Definition~\ref{def:dist-trie}.
\end{definition}

% ------------------------------------------------------------------
\subsection{Alignment}
% ------------------------------------------------------------------

\begin{definition}[Alignment]\label{def:alignment}
  Given a located trace $\trace$ and a trie model~$\T$, an
  \emph{alignment}
  $\alignment = \langle m_1, \ldots, m_p \rangle$
  is a sequence of \emph{moves}
  $(x_j^{\mathrm{model}},\, x_j^{\mathrm{log}})$
  of three types ($\ell = (a, n_i) \in \A \times \Nodes$):
  \begin{alignat*}{3}
    &\textbf{Synchronous:}\quad
    &&\bigl(\ell,\;\ell\bigr),
    &\quad \cost &= 0,\\
    &\textbf{Model move:}\quad
    &&\bigl(\ell,\;\SKIP\bigr),
    &\quad \cost &= 1,\\
    &\textbf{Log move:}\quad
    &&\bigl(\SKIP,\;\ell\bigr),
    &\quad \cost &= 1.
  \end{alignat*}
  The alignment cost is
  $\cost(\alignment) = \sum_{j=1}^{p} \cost(m_j)$.
  An alignment is \emph{optimal} if no other alignment for the same
  $(\trace, \T)$ has strictly lower cost.
\end{definition}

\textbf{Assumption (unit-cost).}\;
We assume
$\cost(\text{model move}) = \cost(\text{log move}) = 1$
throughout.

\begin{definition}[Alignment construction]\label{def:alignment-construction}
  Given an alignment $\alignment$ and a new move~$m$, the extended
  alignment $\alignment' = \alignment \cdot m$ is created by appending
  $m$ to the move sequence.  The cost is tracked incrementally:
  $\cost(\alignment') = \cost(\alignment) + \cost(m)$.
\end{definition}

\begin{oldbox}
\textbf{Flaw: Deep-copy construction.}\;
When using \texttt{deepcopy} on the entire
alignment for every new move, for a trace of length~$m$ with branching
factor~$b$ in A$^{*}$, this produced $O(b^m \cdot m)$ copied elements,
causing exponential memory consumption.
\end{oldbox}

\begin{newboxelem}
\textbf{Fix: Shallow-copy extension.}\;
Since alignment elements are immutable value objects, extending an
alignment copies only the element \emph{list} (not the elements
themselves) and appends the new move.  Cost is tracked as an integer
field, not recomputed.  This reduces per-move cost from $O(m)$ deepcopy
to $O(1)$ amortised append.
\end{newboxelem}

\begin{definition}[Covered events]\label{def:covered}
  Given an alignment $\alignment = \langle m_1, \ldots, m_p \rangle$,
  the set of \emph{covered events} is
  \[
    \cov(\alignment) = \{ x_j^{\mathrm{model}} \mid x_j^{\mathrm{model}} \neq \SKIP \}
    \;\cup\;
    \{ x_j^{\mathrm{log}} \mid x_j^{\mathrm{log}} \neq \SKIP \}.
  \]
\end{definition}

\begin{oldbox}
\textbf{Flaw: One-sided coverage check.}\;
The original \texttt{append\_missing\_log\_moves} only checked the
\emph{log side} of existing alignment elements:
$\cov_{\text{old}} = \{ x_j^{\mathrm{log}} \mid x_j^{\mathrm{log}} \neq \SKIP \}$.
This missed events already handled as synchronous or model moves, causing
duplicate log moves in the final alignment.
\end{oldbox}

\begin{newboxelem}
\textbf{Fix: Two-sided coverage.}\;
$\cov(\alignment)$ includes events from \emph{both} the model and log
sides of every move (Definition~\ref{def:covered}).  An external event
is added as a log move only if it does not appear in $\cov(\alignment)$.
\end{newboxelem}

\begin{definition}[Central alignment]\label{def:central-alignment}
  Let $\Tc$ be a centralised trie built from training traces
  $L_{\mathrm{train}}$.
  The \emph{central optimal alignment} of a trace~$\trace$ is
  \[
    \alignment^{*}_c(\trace)
    \;=\;
    \arg\min_{\alignment \in \Gamma(\trace,\,\Tc)}
      \cost(\alignment),
  \]
  where $\Gamma(\trace, \Tc)$ is the set of all valid alignments
  of $\trace$ w.r.t.\ $\Tc$, computed via A$^{*}$ over the
  product space (trie state $\times$ trace index).
\end{definition}

%% =====================================================================
\section{Formal Problem Statement}
\label{sec:problem}
%% =====================================================================

\begin{problem}[Distributed online conformance checking]\label{prob:dcc}
  Given training traces $L_{\mathrm{train}}$ over~$\A$ and
  $\Nodes = \{n_1, \ldots, n_k\}$, construct a network
  $\Net = (\Nodes, \{\Tl_{n_i}\})$
  such that for every validation trace~$\trace$ and every prefix
  $\trace_{1..i}$, upon receiving the $i$-th event at node~$n_{j_i}$,
  the system produces a distributed alignment $\alignment_i$ with
  \[
    \cost(\alignment_i) = \cost\bigl(\alignment^{*}_c(\trace_{1..i})\bigr),
  \]
  where $\alignment^{*}_c$ is the optimal central alignment
  (Definition~\ref{def:central-alignment}) using a centralised trie
  $\Tc$ built from the same~$L_{\mathrm{train}}$.

  The computation must be \emph{distributed}: each node~$n_i$ sees only
  its own local events $\trace|_{n_i}$ and communicates with other nodes
  solely through the \GA{} protocol (Definition~\ref{def:get-alignment}).
\end{problem}

\medskip\noindent
The core requirement is \emph{cost equivalence}: the distributed
alignment must have the same cost as the centralised one.
This demands that the decomposition into local tries preserves the
structure of the global model, and that the composition of local
alignments yields a globally optimal solution.

\begin{definition}[Trace reconstruction]\label{def:reconstruction}
  The \emph{reconstruction operator}
  $\rho\colon \prod_{n_i} \Lg(\Tl_{n_i}) \to (\A \times \Nodes)^{*}$
  composes local fragments into full located traces by stitching at
  entry points: if $\Tl_{n_i}$ has entry point~$e$ with
  $\lambda(e) = \na{n_j}{a'}$, the local fragment at~$n_i$ is
  concatenated after the fragment at~$n_j$ ending with~$a'$, respecting
  temporal order.
\end{definition}

\noindent
The decomposition must satisfy the language-preservation property:
\begin{equation}\label{eq:lang-preservation}
  \Lg(\Tc)
  = \rho\Bigl(
      \textstyle\prod_{n_i \in \Nodes} \Lg(\Tl_{n_i})
    \Bigr).
\end{equation}

%% =====================================================================
\section{Mechanism: The \GA{} Protocol}
\label{sec:mechanism}
%% =====================================================================

We now describe the distributed alignment protocol that solves
Problem~\ref{prob:dcc}.

% ------------------------------------------------------------------
\subsection{Protocol Definition}
% ------------------------------------------------------------------

\begin{definition}[\GA{} protocol]\label{def:get-alignment}
  Each node $n_\ell$ exposes
  $\GA(n_\ell, e, t, V)$
  returning an \texttt{AlignmentResponse}
  $(\alignment_{\mathrm{up}}, t', e, n_{\mathrm{last}})$ where:
  \begin{itemize}
    \item $e \in \ep(n_\ell)$ is an entry-point label with
          $\loc(e) = n_\ell$,
    \item $t$ is the logical timestamp upper bound,
    \item $V \subseteq \Nodes$ is the \emph{visiting set}
          (nodes already on the current call stack),
    \item $\alignment_{\mathrm{up}}$ is the best alignment of
          $\trace|_{n_\ell}^{< t}$ against the subtrie reaching~$e$,
    \item $t'$ is the timestamp of the last event consumed,
    \item $n_{\mathrm{last}}$ is the node that produced the deepest
          log moves.
  \end{itemize}
\end{definition}

\begin{definition}[Cycle detection via visiting set]%
  \label{def:cycle-detection}
  Before node~$n_\ell$ issues a recursive call
  $\GA(n_k, e', t, V)$, it checks $n_k \notin V$.
  If $n_k \in V$, the entry point is skipped (the candidate set
  excludes this branch).  Each call extends the set:
  $V' = V \cup \{n_\ell\}$.
\end{definition}

\begin{oldbox}
\textbf{Flaw: No cycle detection.}\;
The original protocol had no visiting set.  If node~$A$ had an entry
point labelled with a location at~$B$, and $B$ had an entry point
labelled with a location at~$A$, the mutual recursion
$A \to B \to A \to B \to \cdots$ never terminated, producing a stack
overflow.
\end{oldbox}

\begin{newboxelem}
\textbf{Fix.}\;
A set $V$ is threaded through every recursive \GA{} call.
A node $n$ is added to~$V$ before recursing; any call to a node
already in~$V$ is pruned.  This guarantees each node appears at most
once on any call path.
\end{newboxelem}

% ------------------------------------------------------------------
\subsection{Caching}
% ------------------------------------------------------------------

\begin{definition}[Response cache]\label{def:cache}
  Each node $n_i$ maintains a cache
  $C_{n_i}\colon (\A \times \Nodes) \times \mathbb{N}
    \to \texttt{AlignmentResponse}$
  mapping (entry-point label, timestamp upper bound) pairs to previously
  computed responses.

  \new{The cache is cleared at the start of each \texttt{process\_event}
  call}, ensuring no stale upstream results persist across events.
\end{definition}

\begin{oldbox}
\textbf{Flaw: Stale cache.}\;
The original cache was keyed only by entry-point label, not by
timestamp.  Once a response was cached for entry point~$e$, all
subsequent events—at arbitrarily later times—reused the same stale
response, ignoring new upstream events.
\end{oldbox}

\begin{newboxelem}
\textbf{Fix.}\;
The cache key is the pair $(e, t)$.  Additionally, the entire cache
is invalidated at the start of each new event processing, so upstream
results always reflect the current state of the network.
\end{newboxelem}

% ------------------------------------------------------------------
\subsection{Distributed Online Alignment Procedure}
% ------------------------------------------------------------------

\begin{definition}[Distributed online alignment]%
  \label{def:dist-online-align}
  Given $\Net = (\Nodes, \{\Tl_{n_i}\})$ and a located trace
  $\trace = \langle \ell_1^{t_1}, \ldots, \ell_m^{t_m} \rangle$,
  upon receiving event $\ell_i = (a_i, n_{j_i})$ at time $t_i$,
  node~$n_{j_i}$ computes alignment $\alignment_i$ for
  $\trace_{1..i}$ as follows:

  \begin{enumerate}
    \item \textbf{Entry-point enumeration.}\;
      Identify all children $e$ of the root of $\Tl_{n_{j_i}}$ whose
      subtrie contains $\ell_i$ (i.e.\ $\ell_i \in R(e)$).

    \item \textbf{Upstream alignment.}\;
      For each entry point $e$ with $\loc(e) = n_\ell \neq n_{j_i}$,
      request $\GA(n_\ell, e, t_i, \{n_{j_i}\})$.
      If $n_\ell \in V$, skip this entry point (cycle avoidance).

    \item \textbf{Local A$^{*}$ alignment.}\;
      Compute a local alignment of the \new{full local trace}
      $\trace|_{n_{j_i}}^{< t_i}$ against the subtrie below~$e$,
      targeting label~$\ell_i$.

    \item \textbf{Composition.}\;
      Concatenate upstream and local alignments:
      $\alignment^{(e)} = \alignment_{\mathrm{up}}(e) \cdot
       \alignment_{\mathrm{loc}}(e)$.

    \item \textbf{External log moves} (top level only).
      \new{Only the initiating node} — the node that received the
      original event — adds log moves for uncovered events at other nodes:
      \[
        \alignment'^{(e)} = \alignment^{(e)} \cdot
        \bigl\langle (\SKIP, \ell) \mid
          \ell \in \textstyle\bigcup_{n \neq n_{j_i}} \trace|_n^{< t_i}
          \setminus \cov(\alignment^{(e)}) \bigr\rangle.
      \]

    \item \textbf{Selection.}\;
      Select the minimum-cost candidate over all entry points:
      \[
        \alignment_i
        = \arg\min_{e \in \mathrm{children}(q_0^{n_{j_i}})}
          \cost(\alignment'^{(e)}).
      \]
  \end{enumerate}
\end{definition}

\begin{oldbox}
\textbf{Flaw 1: External log moves added at every recursion level.}\;
The implementation called \texttt{\_add\_external\_log\_moves}
inside every recursive \GA{} call, not just at the top level.  If
node~$A$ added log moves for $B$'s events, and $B$ added log moves
for $A$'s events, the composed alignment contained duplicate log moves,
inflating the cost.
\end{oldbox}

\begin{newboxelem}
\textbf{Fix.}\;
External log moves are added \emph{only} in \texttt{process\_event}
(the top-level entry point).  Recursive \GA{} calls return the
upstream alignment without log-move patching.
This ensures each external event is counted exactly once.
\end{newboxelem}

\begin{oldbox}
\textbf{Flaw 2: Local trace filtered by upstream timestamp.}\;
The original local trace retrieval was
$\trace|_{n_i}^{(t_{\mathrm{up}},\, t_i)}$—only events
\emph{after} the upstream's last timestamp.  If local events occurred
before the upstream transition, they were silently dropped.
\end{oldbox}

\begin{newboxelem}
\textbf{Fix.}\;
The local trace includes \emph{all} events at the current node up to
time~$t_i$, regardless of the upstream timestamp:
$\trace|_{n_i}^{(-\infty,\, t_i)}$.
The upstream timestamp informs only when the upstream alignment ends,
not which local events to consider.
\end{newboxelem}


%% =====================================================================
\section{Proof of Soundness}
\label{sec:proof}
%% =====================================================================

We prove that the \GA{} protocol produces optimal alignments
via a series of lemmas.

% ------------------------------------------------------------------
\subsection{Termination}
% ------------------------------------------------------------------

\begin{lemma}[Termination]\label{lem:termination}
  For any event $\ell_i^{t_i}$ arriving at node~$n_{j_i}$, the
  recursive \GA{} invocations terminate in at most $|\Nodes|$ steps.
\end{lemma}

\begin{proof}
  The visiting set $V$ is initialised as $\{n_{j_i}\}$ and
  grows by one element at each recursive call
  (Definition~\ref{def:cycle-detection}).
  A recursive call $\GA(n_k, e', t, V)$ is issued only if
  $n_k \notin V$.
  Since $|V| \leq |\Nodes|$ and $V$ grows monotonically,
  the recursion depth is bounded by~$|\Nodes| - 1$.
  At the base case, a node with only local starts ($\ep(n) = \emptyset$)
  or all predecessors already in~$V$ computes a purely local A$^{*}$
  alignment without further recursion.
\end{proof}

% ------------------------------------------------------------------
\subsection{Cost Additivity}
% ------------------------------------------------------------------

\begin{lemma}[Cost additivity under composition]\label{lem:additivity}
  Let $\alignment_1$ and $\alignment_2$ be alignments.  Under unit costs,
  \[
    \cost(\alignment_1 \cdot \alignment_2)
    = \cost(\alignment_1) + \cost(\alignment_2).
  \]
\end{lemma}

\begin{proof}
  Each move's cost depends only on its type (synchronous $= 0$,
  model/log $= 1$).  Concatenation appends the move sequences:
  $\alignment_1 \cdot \alignment_2 =
    \langle m_1^{(1)}, \ldots, m_p^{(1)},
            m_1^{(2)}, \ldots, m_q^{(2)} \rangle$.
  The total cost is
  $\sum_{j=1}^{p} \cost(m_j^{(1)}) + \sum_{j=1}^{q} \cost(m_j^{(2)})
  = \cost(\alignment_1) + \cost(\alignment_2)$.
\end{proof}

% ------------------------------------------------------------------
\subsection{Exhaustive Enumeration}
% ------------------------------------------------------------------

\begin{lemma}[Entry-point exhaustiveness]\label{lem:exhaustive}
  Let $\trace^{*}$ be a path in $\Tc$ that passes through node~$n_i$.
  Then there exists exactly one entry point
  $e \in \ep(n_i) \cup \ls(n_i)$ such that the local fragment
  of~$\trace^{*}$ at~$n_i$ is a path in the subtrie below~$e$
  in~$\Tl_{n_i}$.
\end{lemma}

\begin{proof}
  By construction of the distributed decomposition
  (Definition~\ref{def:dist-trie}):
  \begin{itemize}
    \item If $\trace^{*}$ starts at $n_i$, the first activity is a
          local start in $\ls(n_i)$.
    \item If $\trace^{*}$ transitions from $n_j$ to $n_i$ via activity
          $\na{n_j}{a'}$, then $\na{n_j}{a'}$ is an entry point in
          $\ep(n_i)$, and the continuation at~$n_i$ is the subtrie
          below this entry point.
  \end{itemize}
  Uniqueness: each transition point in a path determines exactly one
  entry point (or local start) in the local trie.
  Completeness: every transition from a path in $\Tc$ is represented
  as an entry point or local start in the corresponding local trie.
\end{proof}

% ------------------------------------------------------------------
\subsection{Local Optimality}
% ------------------------------------------------------------------

\begin{lemma}[A$^{*}$ local optimality]\label{lem:local-opt}
  Given a local trace $\trace|_{n_i}$ and a subtrie $\Tl'_{n_i}$
  rooted below an entry point, the A$^{*}$ search over the product
  space $Q_{n_i} \times \{0, \ldots, |\trace|_{n_i}|\}$
  returns an optimal local alignment.
\end{lemma}

\begin{proof}
  A$^{*}$ with an admissible heuristic (here $h = 0$, i.e.\ uniform-cost
  search) is complete and optimal on finite graphs.
  The product space is finite ($|Q_{n_i}| \cdot (|\trace|_{n_i}|+1)$
  states), and the unit-cost function is non-negative.
  Therefore the first solution popped from the priority queue is optimal.

  \new{The goal check uses full located-activity equality
  (Definition~\ref{def:la-equality}), ensuring the search terminates
  at the correct target node in the trie.}
\end{proof}

\begin{oldbox}
\textbf{Flaw: Goal check ignored location.}\;
The original A$^{*}$ goal check compared only the activity name, not
the full located activity.  This could cause the search to terminate
at a trie node with the right activity but wrong location, producing
a spurious alignment.
\end{oldbox}

% ------------------------------------------------------------------
\subsection{Log-Move Completeness}
% ------------------------------------------------------------------

\begin{lemma}[External log-move completeness]\label{lem:log-moves}
  After step~5 of Definition~\ref{def:dist-online-align},
  every event in $\trace_{1..i}$ is covered by exactly one move
  in~$\alignment'_i$.
\end{lemma}

\begin{proof}
  Partition the events in $\trace_{1..i}$ by their location:
  \begin{itemize}
    \item Events at $n_{j_i}$ (the current node): handled by the local
          A$^{*}$ alignment as synchronous or log moves.
    \item Events at upstream nodes along the chosen entry-point path:
          handled by the upstream \GA{} alignment.
    \item Events at all other nodes: added as log moves in step~5,
          provided they are not already in $\cov(\alignment^{(e)})$.
  \end{itemize}

  The coverage set $\cov$ (Definition~\ref{def:covered}) is
  constructed from \emph{both} sides (model and log) of every existing
  element.  Therefore, events already present as synchronous moves
  (appearing on both sides) or model moves (appearing on the model side)
  are not duplicated.
  Events not in $\cov$ are added exactly once.

  \new{Step~5 is executed only at the top level
  (\texttt{process\_event}), preventing nested addition across
  recursion levels.}
\end{proof}

% ------------------------------------------------------------------
\subsection{Composition Optimality}
% ------------------------------------------------------------------

\begin{lemma}[Alignment composition optimality]\label{lem:composition}
  Let $e$ be an entry point in $\Tl_{n_i}$ with $\loc(e) = n_\ell$.
  Let $\alignment_{\mathrm{up}}^{*}$ be optimal from
  $\GA(n_\ell, e, t, V)$ and $\alignment_{\mathrm{loc}}^{*}$ optimal
  locally below~$e$ (Lemma~\ref{lem:local-opt}).
  If the candidate set enumerates all
  $e \in \ep(n_i) \cup \ls(n_i)$ (Lemma~\ref{lem:exhaustive}), then
  \[
    \alignment_i
    = \arg\min_{e}\;
      \bigl(
        \cost(\alignment_{\mathrm{up}}^{*}(e))
        + \cost(\alignment_{\mathrm{loc}}^{*}(e))
      \bigr)
  \]
  is an optimal alignment for the events decomposable through~$n_i$.
\end{lemma}

\begin{proof}
  By Lemma~\ref{lem:exhaustive}, every decomposition of a global path
  through~$n_i$ corresponds to some entry point
  $e \in \ep(n_i) \cup \ls(n_i)$.
  For each~$e$:
  \begin{itemize}
    \item The upstream alignment $\alignment_{\mathrm{up}}^{*}(e)$ is
          optimal by the inductive hypothesis (or by
          Lemma~\ref{lem:local-opt} at the base case).
    \item The local alignment $\alignment_{\mathrm{loc}}^{*}(e)$ is
          optimal by Lemma~\ref{lem:local-opt}.
    \item By Lemma~\ref{lem:additivity},
          $\cost(\alignment_{\mathrm{up}}^{*} \cdot
                 \alignment_{\mathrm{loc}}^{*})
          = \cost(\alignment_{\mathrm{up}}^{*})
          + \cost(\alignment_{\mathrm{loc}}^{*})$.
  \end{itemize}
  Minimising over all entry points considers every possible
  decomposition, so the minimum is globally optimal among all
  decompositions through~$n_i$.
\end{proof}

% ------------------------------------------------------------------
\subsection{Main Theorem}
% ------------------------------------------------------------------

\begin{theorem}[Soundness of distributed online alignment]%
  \label{thm:soundness}
  Under the unit-cost assumption and with the corrected protocol
  (Definitions~\ref{def:get-alignment}--\ref{def:dist-online-align}),
  for every validation trace $\trace$ and every prefix $\trace_{1..i}$,
  the distributed alignment $\alignment_i$ satisfies
  \[
    \cost(\alignment_i) \leq \cost(\alignment^{*}_c(\trace_{1..i})).
  \]
\end{theorem}

\begin{proof}
  By induction on the recursion depth of \GA{}.

  \medskip\noindent
  \emph{Base case} (depth $0$).\;
  Node~$n_i$ has only local starts ($\ep(n_i) = \emptyset$) or all
  predecessor nodes are in~$V$.
  The alignment is computed by A$^{*}$ on the local product space,
  which is optimal by Lemma~\ref{lem:local-opt}.

  \medskip\noindent
  \emph{Inductive step} (depth $d$).\;
  Assume \GA{} returns optimal sub-alignments at depth $\leq d-1$.
  At depth~$d$, node~$n_i$:
  \begin{enumerate}
    \item Enumerates all entry points — exhaustive by
          Lemma~\ref{lem:exhaustive}.
    \item For each entry point, obtains an optimal upstream alignment
          (by the inductive hypothesis).
    \item Computes an optimal local alignment
          (Lemma~\ref{lem:local-opt}).
    \item Composes them with additive cost
          (Lemma~\ref{lem:additivity}).
    \item Adds external log moves exactly once, covering all remaining
          events (Lemma~\ref{lem:log-moves}).
    \item Selects the minimum over all candidates
          (Lemma~\ref{lem:composition}).
  \end{enumerate}

  The recursion terminates by Lemma~\ref{lem:termination}.
  The resulting $\alignment_i$ has cost no greater than any
  centralised alignment, since every decomposition of a centralised
  path is represented in the candidate set.
\end{proof}

\begin{corollary}
  The protocol terminates in at most $|\Nodes|$ recursive steps per
  event, and each invocation explores a local state space of size
  $|Q_{n_i}| \cdot (|\trace|_{n_i}| + 1)$, which is substantially
  smaller than the centralised state space
  $|Q_c| \cdot (m + 1)$ when activities are distributed across nodes.
\end{corollary}

%% =====================================================================
\section{State Space Complexity}
\label{sec:state-space}
%% =====================================================================

\begin{definition}[Alignment state space]\label{def:state-space}
  A$^{*}$ alignment of a length-$m$ trace against a trie with $|Q|$
  states explores the product space
  $S = Q \times \{0, 1, \ldots, m\}$,
  of size $|S| = |Q| \cdot (m+1)$.
\end{definition}

\begin{proposition}[Centralised state space]\label{prop:central-space}
  For the centralised trie~$\Tc$, in the worst case
  $|Q_c| = O\!\bigl(\sum_{\trace \in L_{\mathrm{train}}} |\trace|\bigr)$
  (no shared prefixes), giving search space
  \[
    |S_c|
    = |Q_c| \cdot (m+1)
    = O\!\Bigl(
        \Bigl(\textstyle\sum_{\trace \in L_{\mathrm{train}}} |\trace|\Bigr)
        \cdot m
      \Bigr).
  \]
\end{proposition}

\begin{proposition}[Distributed state space]\label{prop:dist-space}
  Each node~$n_i$ searches
  $S_{n_i} = Q_{n_i} \times \{0, \ldots, |\trace|_{n_i}|\}$.
  The total effort is bounded by
  \[
    \textstyle\sum_{n_i \in \Nodes}
      |Q_{n_i}| \cdot (|\trace|_{n_i}| + 1).
  \]
  Since both trie states and trace are partitioned,
  $|Q_{n_i}| \ll |Q_c|$ and $|\trace|_{n_i}| \leq m$, so each local
  product is substantially smaller than $|Q_c| \cdot m$.
\end{proposition}

\medskip\noindent
\textbf{Cascading overhead.}\;
The worst case produces a chain of $|\Nodes|$ sequential \GA{} calls,
each blocking on its predecessor, yielding $O(|\Nodes|)$ sequential
A$^{*}$ invocations.
Caching (Definition~\ref{def:cache}) reduces repeated requests to
amortised~$O(1)$.

\medskip\noindent
\textbf{Entry-point enumeration.}\;
If $n_i$ has $E_i = |\ep(n_i)|$ entry points, the worst-case per-event
effort is
$O\bigl(E_i \cdot (|Q_{n_i}| \cdot |\trace|_{n_i}|
  + C_{\mathrm{up}})\bigr)$,
where $C_{\mathrm{up}}$ is the upstream cascade cost.

\medskip\noindent
In summary, the decomposition trades the monolithic space
$|Q_c| \cdot m$ for a sum of smaller local spaces at the cost of
inter-node communication.  It is most effective when
(i) $|\Nodes|$ is moderate,
(ii) activities are well-distributed ($|Q_{n_i}| \ll |Q_c|$), and
(iii) the number of entry points~$E_i$ per node is bounded.
\end{document}